\pagestyle{plain}
\chapter{Introdução}

\section{Saberes e Sistemas Agrícolas Tradicionais}

Ontologicamente, os Saberes e Sistemas Agrícolas Tradicionais (SSAT) transcendem a definição simplista de práticas agrárias pretéritas ou de mera agricultura de subsistência. Sob a ótica da agroecologia e da etnoecologia, estes sistemas configuram-se como complexos ecossistêmicos antropogênicos, forjados pela co-evolução milenar entre grupos humanos e seus biomas. Autores seminais como Toledo e Barrera-Bassols \cite{Toledo2008} conceituam essa interação como a materialização da memória biocultural, uma matriz indissociável onde a diversidade biológica é mantida e amplificada pela diversidade cultural, linguística e simbólica das comunidades. Nesse contexto, o saber não opera como um arquivo estático, mas estrutura-se como um \textit{corpus} dinâmico fundamentado no complexo Conhecimento-Prática-Crença (K-P-B) descrito por Fikret Berkes \cite{Berkes2017}.

Sob esta ótica analítica, o componente de crenças, frequentemente marginalizado por análises tecnocráticas como elemento subjetivo, desempenha uma função estruturante crítica. Ele atua como uma cosmovisão normativa que impõe restrições éticas à apropriação dos recursos naturais, transformando a paisagem em um sujeito de direitos e reciprocidade, e não mero objeto de exploração. A indissociabilidade dessa tríade engendra um sistema cibernético de retroalimentação social e ecológica, onde a observação empírica refinada (\textit{Knowledge}) instrumentaliza as instituições de manejo (\textit{Practice}), as quais são, por sua vez, moduladas e legitimadas por valores culturais (\textit{Belief}). Esse arranjo sistêmico impede a sobre-exploração dos estoques biológicos através de tabus e sanções sociais, assegurando que a inovação técnica permaneça subordinada à integridade dos ciclos ecossistêmicos que sustentam a comunidade.

A racionalidade produtiva que rege esses sistemas difere radicalmente da lógica industrial de simplificação ecológica. Enquanto a agricultura moderna busca a homogeneização e o isolamento de variáveis, os SSAT operam sob a lógica da complexidade, da diversidade funcional e da redundância biológica. Miguel Altieri \cite{Altieri1995} descreve esses agroecossistemas como exemplos de engenharia ecológica vernácula, caracterizados por arranjos policulturais, sistemas agroflorestais e manejo integrado de paisagens que maximizam as interações sinérgicas entre solo, plantas e animais. Essa sofisticação sistêmica permite que tais comunidades mantenham altos níveis de produtividade líquida e estabilidade a longo prazo sem dependência massiva de insumos externos, constituindo verdadeiros refúgios de agrobiodiversidade e bancos genéticos \textit{in situ} vitais para o futuro da agricultura global.

Portanto, compreender esses sistemas exige uma abordagem holística que reconheça a sua multifuncionalidade. Eles não produzem apenas \textit{commodities}, mas geram serviços ecossistêmicos, reproduzem identidades territoriais e sustentam estruturas sociais de cooperação. A "tradicionalidade" aqui não denota atraso tecnológico, mas uma trajetória tecnológica distinta, pautada na observação fenológica rigorosa e na experimentação contínua. É justamente essa riqueza de variáveis e interações não-lineares, definida como a complexidade inerente à memória biocultural, que desafia os métodos reducionistas de análise e justifica a adoção de ferramentas computacionais avançadas capazes de decodificar a inteligência ecológica neles depositada.

A gestão estratégica de ativos intangíveis consolidou-se como diferencial competitivo na economia do conhecimento, exigindo novas abordagens para a valoração e proteção de saberes ancestrais. A fundamentação teórica desta premissa remonta à teoria do crescimento da firma de Penrose \cite{Penrose1959}, que identificou a base de recursos internos e as capacidades gerenciais como os limitadores primários e, simultaneamente, os motores do crescimento organizacional. Transpondo essa lógica para o contexto comunitário, os saberes tradicionais deixam de ser vistos como meros traços culturais para serem compreendidos como recursos estratégicos vitais. Na perspectiva da Resource-Based View (RBV) \cite{Barney1991}, a vantagem competitiva sustentável e a sobrevivência de um grupo derivam de recursos que sejam Valiosos, Raros, Inimitáveis e Não-substituíveis (VRIN). O conhecimento tradicional, quando adequadamente gerido e protegido contra apropriações indébitas, atende rigorosamente a esses critérios, configurando-se como um ativo idiossincrático de alto valor.

Contudo, a análise desses ativos demanda superar a visão estática de estoque de conhecimento. Sob a lente da teoria econômica evolucionária, os Saberes e Sistemas Agrícolas Tradicionais (SSAT) materializam o conceito de novas combinações proposto por Joseph Schumpeter \cite{Schumpeter1934}. Longe de representarem o oposto da modernidade, esses sistemas constituem laboratórios vivos de inovação contínua, onde a persistência secular de práticas agroecológicas em ambientes de restrição de recursos evidencia um vigoroso processo de experimentação empírica. Nessa perspectiva analítica, o acervo de conhecimentos detido pelas comunidades transcende o folclore para se configurar como um conjunto robusto de Capacidades Dinâmicas, conforme a estrutura teórica de Teece, Pisano e Shuen \cite{Teece1997}. O monitoramento fenológico minucioso corresponde ao mecanismo de \textit{sensing} (monitoramento de sinais fracos), a seleção de genótipos adaptados opera como o \textit{seizing} (mobilização de recursos) e a modificação de práticas de cultivo reflete o \textit{transforming} (reconfiguração de ativos).

O desafio central para a ciência contemporânea reside na natureza epistemológica desses ativos. Conforme postulado por Michael Polanyi \cite{Polanyi1966}, grande parte desse conhecimento reside em uma dimensão tácita e inefável, operando sob a lógica de que ``sabemos mais do que somos capazes de expressar''. Diferentemente da inovação formal, codificada em patentes e manuais, a inovação nos sistemas tradicionais é transmitida oralmente e incorporada na prática cotidiana (\textit{learning-by-doing}), o que frequentemente a torna invisível às métricas da ciência hegemônica e aos instrumentos clássicos de propriedade intelectual. A crise ambiental global e os imperativos do Antropoceno, entretanto, revelam que esses modelos de gestão tácita incorporam dimensões socioambientais de resiliência que a agricultura industrial falhou em replicar, tornando urgente o desenvolvimento de mecanismos de ``externalização'', conforme o modelo de Nonaka \cite{Nonaka1995}, capazes de traduzir essa complexidade para linguagens validáveis.

Operacionalizar essa visão de conhecimentos tradicionais como ativos estratégicos e dinâmicos demanda, portanto, instrumentos jurídicos e, primordialmente, metodologias tecnológicas de decodificação. Se o conhecimento tradicional configura-se como um sistema adaptativo complexo e \textit{path-dependent} \cite{Nelson1982}, sua gestão eficaz exige o desenvolvimento de capacidades absortivas organizacionais \cite{CohenLevinthal1990}, permitindo que as comunidades integrem tecnologias digitais sem comprometer sua autenticidade cultural. O paradigma da Inovação Aberta \cite{Chesbrough2003} fornece o arcabouço para essa integração, onde o fluxo de conhecimento entre a ciência de dados e a sabedoria ancestral não ocorre por imposição, mas por processos de co-criação que visam salvaguardar o patrimônio biocultural.

A valoração de conhecimentos tradicionais como ativos intangíveis de inovação demanda metodologias sofisticadas que transcendam as abordagens econométricas lineares ou puramente baseadas em custos. A literatura sobre valoração de ativos evoluiu para frameworks analíticos que incorporam a complexidade multidimensional da inovação, conforme discutido por Crossan e Apaydin \cite{Crossan2010}. Nesse cenário, a densidade de interações biológicas, climáticas e sociais presentes nos sistemas tradicionais excede a capacidade de apreensão por métodos de observação estatística convencional, justificando imperativamente a aplicação de ferramentas de Engenharia de Dados e Inteligência Artificial.

O Aprendizado de Máquina (Machine Learning) emerge neste contexto não apenas como uma ferramenta instrumental de processamento, mas como um`Decodificador Tecnológico de Capacidades Dinâmicas. Ao processar a alta dimensionalidade de dados fenológicos e ambientais, algoritmos avançados possuem a capacidade hermenêutica de identificar padrões latentes que constituem os microfundamentos da adaptação tradicional. A aplicação dessas tecnologias permite explicitar o conhecimento tácito, transformando observações empíricas dispersas em parâmetros científicos auditáveis e validáveis. Esse processo de tradução tecnológica é fundamental para fundamentar a legitimidade científica dos Saberes e Sistemas Agrícolas Tradicionais, criando a base probatória necessária para sua proteção jurídica e valorização econômica.

A salvaguarda dos saberes dos Povos e Comunidades Tradicionais (PCTs) representa, portanto, um desafio estratégico de gestão da inovação e de justiça territorial. No Brasil, apesar dos avanços constitucionais e da ratificação de tratados como a Convenção 169 da OIT, a insegurança jurídica territorial permanece crítica. A ausência de documentação técnica padronizada que comprove a ocupação ancestral e a gestão sofisticada dos recursos naturais frequentemente fragiliza a posição dessas comunidades em disputas fundiárias e processos de licenciamento. Nesse sentido, o mapeamento sistemático e a modelagem computacional dos conhecimentos agroecológicos não servem apenas à curiosidade acadêmica, mas constituem evidências materiais robustas de ocupação e uso sustentável do território.


A relevância desta investigação reside na proposição de um framework inovador que instrumentaliza a Engenharia de Dados para o empoderamento comunitário e a proteção de ativos estratégicos. A gestão de conhecimentos tradicionais, quando suportada por evidências computacionais robustas, fortalece a governança territorial e a soberania epistemológica das comunidades. O projeto contribui para a consolidação de capacidades territoriais de inovação social, alinhando-se a modelos de hélice tripla que integram universidades, comunidades e instituições de pesquisa em prol do desenvolvimento sustentável.

Nesse contexto, a presente tese estrutura-se sob um delineamento metodológico processual e integrativo. A investigação inicia-se pela mineração de dados e revisão sistemática da literatura, visando diagnosticar o estado da arte na aplicação de inteligência artificial em contextos agroecológicos. Subsequentemente, avança para a modelagem estatística multivariada e análise de correspondência, buscando identificar os perfis de especialização e as capacidades dinâmicas ocultas nas práticas comunitárias. O processo evolui para a validação participativa e co-criação metodológica, garantindo que os modelos computacionais reflitam a realidade fenomenológica dos territórios, e culmina na otimização multiobjetivo para a proposição de mecanismos de salvaguarda que equilibrem preservação cultural e aplicabilidade científica.

Esta abordagem integrada visa, em última instância, superar a lacuna entre a riqueza empírica dos saberes tradicionais e a exigência de formalização da ciência moderna, estabelecendo a inteligência artificial como uma ponte para o reconhecimento de direitos e a valorização de tecnologias sociais ancestrais.


\section{Problema de Pesquisa}

O problema central reside na falha de mercado decorrente da assimetria informacional entre os detentores de ativos tradicionais e os agentes econômicos e regulatórios. Embora os sistemas produtivos quilombolas gerem serviços ecossistêmicos e produtos de alta complexidade (ativos VRIN) \cite{Barney1991}, a inexistência de mecanismos formais de mensuração e sinalização (\textit{signaling}) impede a captura desse valor intangível. Sem ferramentas auditáveis capazes de converter conhecimento tácito e difuso em parâmetros verificáveis, esses ativos permanecem subprecificados, vulneráveis à apropriação indevida e excluídos de mercados de maior valor agregado.

Essa lacuna desdobra-se em fricções institucionais e metodológicas que se retroalimentam. Em primeiro lugar, a assimetria regulatória e os custos de transação associados aos processos de \textit{compliance} criam barreiras de entrada para comunidades com baixa capacidade técnico-institucional. Em segundo lugar, verifica-se o esgotamento de métodos lineares de elicitação para capturar fronteiras difusas e dimensões simbólicas do conhecimento tradicional, demandando modelagens computacionais capazes de operar com ambiguidade e incerteza. Por fim, a ausência de sinais de mercado confiáveis (origem, autenticidade e sustentabilidade) sustenta seleção adversa, reduzindo incentivos à preservação e à governança territorial.


\section{Questões de Pesquisa}

\begin{itemize}

\item \textbf{Q1 – Prontidão Operacional do Aprendizado de Máquina:}
Em que medida as técnicas de aprendizado de máquina apresentam prontidão operacional para monitoramento e avaliação de sistemas agrícolas tradicionais, considerando acurácia, conformidade FAIR e protocolos de validação espacial?

\item \textbf{Q2 – Decodificação e Estruturação de Saberes Tácitos:}
Como métodos estruturados de adaptação transcultural e elicitação (Delphi) contribuem para converter saberes tácitos em variáveis mensuráveis e computacionalmente processáveis, aumentando consenso e comparabilidade entre avaliadores?

\item \textbf{Q3 – Modelagem Fuzzy e Priorização de Salvaguarda:}
Como a lógica fuzzy pode representar a ambiguidade inerente a dimensões culturais e simbólicas, produzindo um índice de prioridade de salvaguarda (Índice de Salvaguarda Biocultural, ISB) com maior aderência à percepção comunitária de vulnerabilidade do que classificações lineares?

\item \textbf{Q4 – Auditoria Computacional e Aceitação Institucional:}
Em que medida a documentação computacional gerada pelo sistema (dossiês ISB, fichas WOCAT-SLM-QBR, relatórios fuzzy) atende aos requisitos de completude, rastreabilidade e conformidade normativa exigidos por órgãos de salvaguarda e proteção (IPHAN, INPI, CGen, FAO-GIAHS)?

\item \textbf{Q5 – Validação Territorial e Adoção Tecnológica:}
Como a incorporação de usuários finais na construção das regras e interfaces do sistema computacional (\textit{participatory design}) afeta a apropriação, adoção e uso continuado da tecnologia de decodificação no território?

\end{itemize}


\section{Hipóteses}

\begin{itemize}

\item \textbf{H1 – Prontidão Operacional do Aprendizado de Máquina:}
A aplicação de aprendizado de máquina a sistemas agrícolas tradicionais, embora alcance alta acurácia reportada, apresenta lacunas críticas em conformidade FAIR, validação espacial e explicabilidade (XAI) que limitam sua prontidão operacional para aplicações regulatórias e de governança territorial.

\item \textbf{H2 – Decodificação e Estruturação de Saberes Tácitos:}
A aplicação combinada de adaptação transcultural (WOCAT-SLM-QBR) e método Delphi sobre variáveis culturais produz um conjunto validado de variáveis computacionalmente processáveis, com índice de consenso estatístico superior ao obtido por métodos de levantamento qualitativo não estruturados.

\item \textbf{H3 – Aderência da Modelagem Fuzzy à Percepção de Vulnerabilidade:}
O Índice de Salvaguarda Biocultural (ISB), gerado via lógica fuzzy a partir de variáveis de risco de erosão intergeracional, complexidade biocultural, singularidade territorial e status de documentação, demonstra maior fidelidade representacional à percepção comunitária de urgência de salvaguarda do que classificações lineares ou ordinais.

\item \textbf{H4 – Eficácia Probatória e Aceitação Institucional:}
A documentação computacional auditável gerada pelo sistema (dossiês ISB, fichas WOCAT-SLM-QBR adaptadas e relatórios de inferência fuzzy) atende aos critérios de completude, rastreabilidade e conformidade normativa exigidos para instruir processos de salvaguarda junto a órgãos regulatórios, com taxa de aceitação institucional superior à de documentação convencional não padronizada.

\item \textbf{H5 – Validação Territorial e Adoção Tecnológica:}
A incorporação dos usuários finais na etapa de design das regras de inferência do sistema computacional (\textit{participatory design}) aumenta a apropriação psicológica da tecnologia de decodificação, resultando em taxa de adoção continuada superior à de ferramentas de gestão exógenas.

\end{itemize}



\section{Objetivos}
\subsection{Objetivo Geral}

Desenvolver e validar uma arquitetura computacional integrada, articulando aprendizado de máquina, lógica fuzzy e modelagem psicométrica, para decodificar, documentar e salvaguardar ativos intangíveis em Saberes e Sistemas Agrícolas Tradicionais de comunidades quilombolas do Território de Identidade Semiárido Nordeste II (Bahia).

\subsection{Objetivos Específicos}

\begin{itemize}

\item Mapear e avaliar criticamente o estado da arte na aplicação de aprendizado de máquina para monitoramento e avaliação de sistemas agrícolas tradicionais, identificando lacunas metodológicas em validação espacial, explicabilidade e governança de dados que condicionam a prontidão operacional dessas tecnologias.

\item Mapear e caracterizar as dimensões de vulnerabilidade e complexidade biocultural (biofísicas, culturais, simbólicas e de transmissão intergeracional) dos sistemas tradicionais, utilizando adaptação transcultural de instrumentos e técnicas de elicitação estruturada (Delphi) para converter saberes tácitos em variáveis computacionalmente processáveis e auditáveis.

\item Desenvolver e implementar um motor de inferência baseado em lógica fuzzy, integrando as variáveis mapeadas para gerar um índice quantitativo de prioridade de salvaguarda (Índice de Salvaguarda Biocultural -- ISB), classificando saberes segundo urgência de documentação e proteção.

\item Avaliar a eficácia probatória e a aceitação institucional da documentação computacional gerada pelo sistema (dossiês ISB, fichas WOCAT-SLM-QBR, relatórios fuzzy), verificando conformidade normativa e rastreabilidade perante órgãos de salvaguarda e proteção de propriedade intelectual (IPHAN, INPI, CGen, FAO-GIAHS).

\item Validar a adoção territorial do sistema computacional de salvaguarda mediante design participativo com comunidades quilombolas, avaliando usabilidade, apropriação e sustentabilidade institucional da tecnologia como instrumento de empoderamento e proteção de ativos bioculturais.


\end{itemize}


\section{Justificativa}

A invisibilidade documental dos Saberes e Sistemas Agrícolas Tradicionais (SSAT) perante órgãos regulatórios configura uma falha de governança que perpetua a vulnerabilidade territorial das comunidades detentoras. Embora esses saberes constituam ativos intangíveis de elevada complexidade biocultural, a inexistência de mecanismos computacionais capazes de traduzir conhecimento tácito em parâmetros auditáveis impede a instrução de processos de salvaguarda patrimonial (IPHAN), registro de propriedade intelectual coletiva (INPI), autorização de acesso ao patrimônio genético (CGen) e candidatura a programas internacionais de preservação (FAO-GIAHS). A ausência de documentação padronizada e rastreável fragiliza a posição dessas comunidades em disputas fundiárias, processos de licenciamento e negociações de repartição de benefícios.

Esta pesquisa justifica-se pela adoção da Teoria dos Conjuntos Difusos, proposta por \citeonline{Zadeh1965}, como fundamento matemático para representar graus de pertencimento e fronteiras nebulosas em variáveis qualitativas de vulnerabilidade e complexidade biocultural. Ao articular lógica fuzzy e modelagem psicométrica, o estudo busca converter julgamentos linguísticos em regras operacionais e produzir um sistema de auditoria e salvaguarda que sustente decisões de proteção, priorização e governança em ambientes regulatórios complexos.

Além do mérito científico, a arquitetura metodológica proposta é desenhada para capilaridade acadêmica e formação de competências. A decomposição do problema em módulos interdependentes permite que diferentes trilhas empíricas e computacionais convirjam para um arcabouço de salvaguarda biocultural e para um artefato replicável de suporte à decisão, com potencial de impacto territorial e institucional.

\section{Delineamento do estudo}

Este estudo configura-se como pesquisa aplicada de métodos mistos, com delineamento sequencial e abordagem orientada à construção de artefato (sistema computacional de salvaguarda e auditoria) e sua validação. O delineamento metodológico está estruturado em cinco fases integradas:

\begin{enumerate}[label=\bfseries(\Roman*)]
  \item \textbf{Revisão de escopo e diagnóstico tecnológico:} revisão sistemática (PRISMA-ScR) da literatura sobre aprendizado de máquina aplicado a sistemas agrícolas tradicionais, com meta-análise de acurácia, avaliação de conformidade FAIR e identificação de lacunas em validação espacial e explicabilidade;

  \item \textbf{Elicitação e estruturação de variáveis:} adaptação transcultural do questionário WOCAT-SLM e aplicação de procedimentos Delphi para converter saberes tácitos em variáveis linguísticas auditáveis e critérios mensuráveis;

  \item \textbf{Modelagem fuzzy e construção do ISB:} desenho do motor de inferência fuzzy (funções de pertinência, regras, defuzzificação) e cálculo do Índice de Salvaguarda Biocultural, classificando saberes segundo urgência de proteção;

  \item \textbf{Auditoria computacional e aceitação institucional:} submissão da documentação gerada pelo sistema a especialistas e órgãos regulatórios (IPHAN, INPI, CGen, FAO-GIAHS) para avaliação de completude, rastreabilidade e conformidade normativa;

  \item \textbf{Design participativo e validação territorial:} co-criação das regras e interfaces do sistema de salvaguarda com usuários finais, avaliação de usabilidade, apropriação e adoção, e consolidação do ecossistema de pesquisa aplicada e formação.
\end{enumerate}



